\documentclass[addpoints,twoside]{exam}
\usepackage{problemset}

\begin{document}

\FontEasyRead

\Heading
{Microeconomics*  Pop Quiz}
{\underline{SHORT} Answers on Chapters 3-4}
{\textbf{Instructions: }Answer each question succinctly, and show your reasoning. Open book, notes, and
colleagues, no generative AI. 11 questions, 110 points, 60-75 minutes.}

\begin{questions}

  \Problem{Explain the three key assumptions about consumer preferences: completeness, transitivity,
  and “more is better.”}

  \Problem{Why can't indifference curves (with different utilities) cross? What assumption would be
  violated if they did?}

  \Problem{A consumer has a fixed income and chooses between goods $X$ and $Y$.
    \begin{itemize}
      \item What does the budget line represent?
      \item What determines its slope, and how would you calculate it?
  \end{itemize}}

  \Problem{Explain the concept of marginal rate of substitution (MRS). Why does MRS
  typically \underline{diminish} as you move downward along an indifference curve?}

  \Problem{Explain the substitution effect and the income effect on consumption of a normal good
  when the price of that good \underline{rises}.}

  \Problem{Define marginal utility. What does it mean for marginal utility to be diminishing?}

  \pagebreak

  \Problem{Define the following types of goods and explain how demand changes with income:
    \begin{itemize}
      \item Normal good
      \item Neutral good
      \item Inferior good
  \end{itemize}}

  \Problem{A consumer is willing to give up 3 units of $Y$ to gain 1 unit of $X$.
    \begin{itemize}
      \item What is the marginal rate of substitution of $X$ for $Y$?
      \item If $\frac{P_X}{P_Y} = 2$, is the consumer at an optimal bundle? Explain.
  \end{itemize}}

  \Problem{What is an Engel curve? What does its slope tell you about a good?}

  \Problem{Suppose Consumer A demands 6 units at a given price and Consumer B demands 4 units. There are no
    other consumers.
    \begin{itemize}
      \item What is total market demand?
      \item What does “horizontal summation” mean in this context?
  \end{itemize}}

  \Problem{Explain why a linear demand curve does not have constant elasticity. (Hint: Refer to the
    formula for point elasticity and explain what happens when P and Q change.) Give an example of
  a function (not linear!) defining a demand curve that \underline{does} have constant elasticity.}

\end{questions}

\end{document}
